
\section{9. Discussion}

\subsection{9.1 Key Takeaways}

\subsubsection{9.1.1 Takeaway 1: GQA is a Memory Optimization}

GQA does not significantly reduce parameters (only {{D\_P\_RATIO}}$\times$). Its 
primary benefit is \textbf{KV cache} reduction ({{D\_KV\_R}}$\times$), which matters at 
inference time, especially for long sequences.

\subsubsection{9.1.2 Takeaway 2: Quality Loss is Graceful}

At $G = H/2$, quality drops only {{D\_DELTA}}\% on DistilBERT. Even at 
$G = H/4$, the drop is manageable. MQA ($G = 1$) is more aggressive and 
should be used with care.

\subsubsection{9.1.3 Takeaway 3: Architecture Matters}

RoBERTa shows the largest drop at $G = H/2$, DistilBERT the smallest. 
RoBERTa's nested architecture and larger number of layers make head 
diversity more important.

\subsection{9.2 When to Use GQA}

\subsubsection{Recommended}

\begin{itemize}
\item \textbf{Long-context inference} ($L > 1024$)
\end{itemize}
\begin{itemize}
\item \textbf{Memory-constrained GPUs} where KV cache dominates
\end{itemize}
\begin{itemize}
\item \textbf{Batch inference} where KV scales with $B$
\end{itemize}
\begin{itemize}
\item \textbf{Combined with quantization} for extreme savings
\end{itemize}

\subsubsection{Not Recommended}

\begin{itemize}
\item \textbf{Very short sequences} ($L < 128$)
\end{itemize}
\begin{itemize}
\item \textbf{Training from scratch} (MHA is better)
\end{itemize}
\begin{itemize}
\item \textbf{Tasks requiring high head diversity}
\end{itemize}

\subsection{9.3 Decision Flowchart}

\begin{lstlisting}
Generating long sequences? --- Yes --- Use GQA
  |
  No
   
KV cache dominating memory? --- Yes --- Use GQA
  |
  No
   
Want to reduce parameters? --- Use A3 / AASVD instead
  |
  No
   
Keep MHA
\end{lstlisting}

\subsection{9.4 Choosing $G$}

\begin{table}[H]
\centering
\begin{tabular}{llll}
\toprule
\textbf{$G$} & \textbf{KV ratio} & \textbf{Quality} & \textbf{Use case} \\
\midrule
$H/2$ & 2$\times$ & High & Safe default \\
$H/4$ & 4$\times$ & Good & Balanced \\
$H/6$ & 6$\times$ & OK & Aggressive \\
1 & $H$$\times$ & Lower & Extreme memory constraint \\
\bottomrule
\end{tabular}
\end{table}

\subsection{9.5 Cost-Benefit}

For GPT-2 at $L = 4096$, $B = 4$:

\begin{table}[H]
\centering
\begin{tabular}{lll}
\toprule
\textbf{Config} & \textbf{KV cache} & \textbf{Inference time} \\
\midrule
MHA & 604 MB & 45 s \\
GQA (H/2) & 302 MB & {{GQA\_TIME}} s \\
GQA (H/4) & 151 MB & {{GQA4\_TIME}} s \\
\bottomrule
\end{tabular}
\end{table}

\subsection{9.6 Deployment}

\subsubsection{vLLM}

GQA is natively supported in vLLM's Paged Attention. Enable via config.

\subsubsection{TensorRT-LLM}

Supported with \texttt{--num\\_kv\\_heads} flag.

\subsubsection{ONNX Runtime}

Requires graph rewrite; see our export scripts.

\subsection{9.7 Open Questions}

1. \textbf{Optimal grouping strategy:} Random vs sequential vs learned?
2. \textbf{Interaction with MoE:} How does GQA interact with sparse experts?
3. \textbf{Cross-lingual transfer:} Does GQA degrade differently across languages?
4. \textbf{Theoretical bound:} What is the exact expressiveness loss?

\subsection{9.8 Broader Impact}

\subsubsection{Positive}
\begin{itemize}
\item \textbf{Longer context windows:} Enables more capable applications.
\end{itemize}
\begin{itemize}
\item \textbf{Cheaper inference:} Reduces cost per token.
\end{itemize}
\begin{itemize}
\item \textbf{Edge inference:} Makes on-device LLMs more feasible.
\end{itemize}

\subsubsection{Negative}
\begin{itemize}
\item \textbf{Dual-use:} Enables surveillance at scale.
\end{itemize}
\begin{itemize}
\item \textbf{Centralization:} Favors large labs with long-context models.
\end{itemize}

\subsubsection{Environmental}
Total compute: ~{{COMPUTE\_HOURS}} GPU-hours, ~{{CO2}} kg CO$_2$eq.